\documentclass[11pt,a4paper]{letter}
\usepackage[margin=1in]{geometry}
\usepackage[T1]{fontenc}
\usepackage[utf8]{inputenc}
\usepackage{lmodern}
\usepackage[hidelinks]{hyperref}

\signature{Jonathan Anthonio Nii Pedro Nelson\\Kwame Nkrumah University of Science and Technology (KNUST)\\Kumasi, Ghana\\\href{mailto:jonathannelson707@gmail.com}{jonathannelson707@gmail.com}\\ORCID: 0009-0006-9393-6237}
\address{Jonathan Anthonio Nii Pedro Nelson\\Kwame Nkrumah University of Science and Technology (KNUST)\\Kumasi, Ghana}
\date{\today}

\begin{document}
\begin{letter}{Prof. Dr. Steven Haberman\\Editor-in-Chief, \textit{Risks}\\MDPI}
\opening{Dear Prof. Dr. Haberman and Editors of \textit{Risks},}

I am pleased to submit the manuscript entitled ``Commodity Tail Dependence and Exchange-Rate Risk in Ghana: Copula Adequacy, Stress Tests, and Calendar Robustness'' for consideration as an Article in \textit{Risks}.

The manuscript examines financial tail risk at the intersection of commodity markets and exchange-rate risk in Ghana. Using daily-source data on cocoa, gold, Brent crude oil, WTI crude oil and the Ghanaian cedi over 2015--2026, the study treats two defensible calendar constructions as co-equal specifications of the observation process. It combines explicit marginal-model adequacy checks, absolute goodness-of-fit testing for eight static bivariate copula families, finite-threshold tail-concentration measures, moving-block-bootstrap stress comparisons, and multiplicity control.

The main contribution is not the selection of a single preferred copula, but the identification of which dependence conclusions remain defensible after calendar construction, marginal adequacy, absolute model fit and multiple testing are considered jointly. The analysis finds no Bonferroni- or Benjamini--Hochberg-adjusted stress--calm discovery under the combined, COVID-19 or 2024 stress definitions in either calendar panel. Gold--Brent, Gold--WTI and Brent--WTI reject all eight tested static copulas under both calendar constructions, while the cocoa-related commodity pairs are calendar-sensitive. Cedi-related copula and extreme-value quantities are reported explicitly as conditional diagnostics because the tested continuous Cedi marginal specifications remain inadequate.

I believe the manuscript fits the scope of \textit{Risks} because it addresses financial risk measurement, exchange-rate risk, tail dependence, econometric model adequacy and robustness in a commodity-dependent emerging-market setting. The paper also emphasizes reproducibility: the audited code, processing scripts and publication-rebuild materials are maintained in the accompanying GitHub repository.

The submission includes a separate supplementary PDF containing three diagnostic figures and an original graphical abstract. The manuscript declares no external funding and no conflicts of interest. The use of OpenAI's ChatGPT in manuscript preparation is disclosed transparently in the Acknowledgments.

\textbf{AUTHOR CONFIRMATION REQUIRED BEFORE UPLOAD:} Please verify the following declaration before this draft is converted to the final cover letter: ``Neither this manuscript nor any part of its substantive content has been published previously or is currently under consideration for publication in another journal.'' If the manuscript has previously been submitted to an MDPI journal, the prior manuscript ID should also be disclosed in the submission system.

Thank you for considering the manuscript for publication in \textit{Risks}.

\closing{Sincerely,}
\end{letter}
\end{document}
